\section{Discussion}\label{sec:discussion}

\paragraph{The manual-intervention mechanism is measured, not explained.}
\S\ref{sec:results} reports that \texttt{full} and \texttt{recovery\_only} increase manual
interventions relative to \texttt{baseline}, the opposite of the replicated paper's claim. We
verified the measurement (the metric is a direct count of
\texttt{telemetry.\allowbreak manual\_interventions} rows, the same table every configuration writes
to) but not its mechanism: whether repeated
escalation across ticks of one open fault, rate-limit or budget denials converting into
human hand-offs, or some interaction between the two, dominates. Distinguishing these would need a
per-tick trace of proposal outcomes across the fault lifetime, which the current telemetry does not
retain at that grain. We record this as open rather than adjudicate it from the aggregate count; a
reader relying on this system's manual-intervention claim should treat it as unresolved, not resolved
in the system's favour.

\subsection{Limitations and threats to validity}\label{sec:limitations}

\paragraph{Threats to validity.}
Four apply directly. \emph{Simulated human.} Every comparison against \texttt{baseline} inherits the
log-normal on-call model (\S\ref{sec:sensitivity} gives its own break-even, not a defence of the
model). \emph{One live model family.} The live arm used one reasoning model and one faster model from
one provider at temperature zero; \S\ref{sec:results}'s decision-quality and cache-poisoning findings
(D-104l) are properties of that pairing and may not generalise to another model, a claim this paper
does not make and the cross-model harness (D-063, \S\ref{sec:related}) exists to test but was not run
this campaign. \emph{Unequal replicate counts.} \texttt{baseline}, \texttt{rule\_based},
\texttt{autoscale} and \texttt{full} have $\NNFull{}$ replicates; the four single-agent ablations have
half that, a cost trade-off made before the campaign (\S\ref{sec:method}) that narrows their
confidence intervals less than the headline comparisons. \emph{Nominal seed matching.} Because the
per-run seed is a function of the configuration name (\S\ref{sec:method}), no two configurations share
fault or human-latency randomness at a matched (scenario, replicate) index; we used the unpaired
rank-sum test as primary for exactly this reason, but a reader should not interpret the retained
signed-rank statistic as evidence from a truly paired design.

\paragraph{One node, one policy pin, one set of scenarios.}
The testbed is a single machine; the resource-contention scenario's realism (real CPU stress) is
also why runs cannot be parallelised, which bounds how large a campaign is practical. The policy
engine is pinned to one release, and \S\ref{sec:lessons} (L6) reports a same-vendor newer release
that would have changed the adversarial result had we not pinned; this is disclosed as a property of
this evaluation, not investigated as a policy-engine defect. Four fault types were injected; the
production incident taxonomy a real pipeline sees is larger, and L4 (\S\ref{sec:lessons}) is itself
evidence that the mapping from injected fault type to real incident type is not free.

\paragraph{Ethics and autonomy.}
The decision-quality result most relevant to deployment is not the headline MTTR number but
$\NDeltaDecisionCorrectFull{}$: at these settings, \texttt{full} executes an accepted mitigation in
under half its runs. Combined with the manual-intervention finding above, this campaign gives no
support for running this class of system in an unattended, fully autonomous mode against real
infrastructure; it is evidence \emph{for} the graduated trust ladder --- shadow, approval, autonomous,
with a durable kill switch and a per-target action cap (Appendix~\ref{app:deviations}, D-065) --- that
this project already treats as the safe production default, and against the research benchmark's
default of full autonomy ever being read as a deployment recommendation.

\paragraph{What would change our mind.}
A cross-model repetition of \S\ref{sec:livemock} that found the same order-of-magnitude live/mock gap
with a different provider would strengthen the paper's central claim considerably; one that did not
would be evidence the gap is specific to this model pairing rather than to mock evaluation as a
method, and would need reporting just as plainly as the current result.
